\par
Von den in der letzten Briefstelle erwähnten beiden Abhandlungen interessiert
uns hier besonders die zweite, die unter dem Titel \textit{Summatio quarumdam
serierum singularium} im August 1808 der Gesellschaft der Wissenschaften vorgelegt
worden ist\footnote{Erschienen 1811 im Bande I der Commentat. rec., Werke II, S. 9, \textit{Anzeige} ebenda S. 155.}. Nach der Tagebuchnotiz 118. von Mitte Mai 1801 hat
\textsc{Gauss} damals den den Kern der Untersuchung bildenden Satz vermutet, aber
die Tagebuchnotiz 123. vom 30. August 1805 (die erste nach der dreijährigen
Unterbrechung) besagt: »Der Beweis des wunderschönen, oben im Mai 1801
erwähnten Lehrsatzes, den wir während mehr als vier Jahren mit aller Anstrengung
gesucht hatten, ist uns endlich gelungen«\footnote{Vergl. den Brief an \textsc{Olbers} vom 3. September 1805, Werke X, 1, S. 24. Siehe auch \textsc{Bachmann}, Abh. 1 dieses Bandes, Art. 18, S. 45.}. Es handelt sich dabei
um die Bestimmung des Vorzeichens bei der Auswertung der sogenannten
\textsc{Gaussschen} Summen; \textsc{Gauss} dürfte wohl von vornherein einen Zusammenhang
zwischen dieser Vorzeichenbestimmung und der Bestimmung der achten Einheitswurzel
(siehe Werke X, 1, S. 225, Art. 8 das \(\varepsilon=\sqrt{i}\), vergl. oben S. 99),
die bei der linearen Transformation der summatorischen Funktionen \(p(t),\ q(t),\)
\(r(t)\) auftritt, vermutet haben. Die vier Jahre hindurch fortgesetzten Bemühungen,
die \textsc{Gauss} der Bestimmung jenes Vorzeichens gewidmet hat, haben ihn sicherlich
veranlasst, alle ihm zu Gebote stehenden Hilfsmittel heranzuziehen, und
wenn man bedenkt, dass \textsc{Hermite}\footnote{\textsc{Liouvilles} Journal (2), t. 3, \textsc{Hermite} Oeuvres I, S. 486.} die Bestimmung der achten Einheitswurzel
bei der linearen Transformation gerade auf die \textsc{Gaussschen} Summen
zurückgeführt hat, so würde man es kaum begreifen können, dass \textsc{Gauss} bei
den hier in Rede stehenden Fragen seine Untersuchungen über die Funktionen\par
\clearpage\SourceLeaf{208}\ArithmeticHead{123}{R}{VIa. ZAHLENTHEORETISCHE UNTERSUCHUNGEN.}
\noindent \(p,q,r\) nicht herangezogen haben sollte, zumal die Identitäten zwischen Reihen
und Produkten, wie sie in den Artikeln 8, 9 der Abhandlung \textit{Summatio quarumdam
serierum singularium} auftreten, direkt an die Methoden anknüpfen, mit
denen in der Scheda Ac operiert wurde. Die Formel
\stepcounter{footnote}
\[
1+x+x^3+x^6+\text{etc.}=\frac{(1-x^2)(1-x^4)(1-x^6)\ldots}{(1-x)(1-x^3)(1-x^5)\ldots}\text{\textsuperscript{172)}}
\]
\footnotetext{Werke II, S. 20; \textsc{Jacobi} hat diese Formel in Art. 66 der \textit{Fundamenta} hergeleitet und zitiert daselbst die \textit{Summatio quarumdam serierum}, vergl. oben S. 17, Fussnote \textsuperscript{29)}.}
findet sich an mehreren Stellen des Nachlasses, früheren und späteren\footnote{So Scheda Ac (Mai 1800) Werke X, 1, S. 201, Art. 9, Gleichung 3), Handbuch 16, Bb (1808), ebenda S. 292, Art. 3, Gl. (2); Ansätze schon in den \textit{Exercitationes} Art. 9 (1796), ebenda S. 142, und in der Scheda Aa (1798), wo es heisst (vergl. Werke X, 1, S. 277): »Untersuchung der Faktoren der unendlichen Reihe
\[
1+x+x^3+x^6+x^{10}+x^{15}+\cdot\cdot\cdot=S.\text{«}
\]
}, und
die an dieselbe geknüpfte Bemerkung »diese Gleichheit zwischen zwei nicht
auf der Hand liegenden Ausdrücken, auf deren Untersuchung wir bei anderer
Gelegenheit zurückkommen werden, ist doch wohl sehr bemerkenswert« kann
sich auf nichts anderes beziehen, als auf das von \textsc{Gauss} projektierte Werk
über die »transzendenten Funktionen«. Ferner haben wir, allerdings aus
späterer Zeit\footnote{\textit{Hundert Theoreme über die neuen Transzendenten}, Werke III, S. 461 ff. Diese Abhandlung stammt jedenfalls aus der Zeit nach 1818, weil darin (auf S. 467) die \textit{Determinatio attractionis} zitiert wird, wahrscheinlich ist sie um 1825 geschrieben.} herrührend, eine die elliptischen Funktionen betreffende Aufzeichnung
von \textsc{Gauss}, in der direkt auf die \textit{Summatio quarumdam serierum} Bezug
genommen wird. \textsc{Gauss} hat also den Zusammenhang zwischen den Reihen,
die in dieser Abhandlung auftreten, und den elliptischen Funktionen sicher
gekannt.

In dem im November 1801 begonnenen Handbuche (16, Bb) finden sich
Aufzeichnungen zur Theorie der lemniskatischen Funktionen, die Werke III,
S. 405, Art. [4] und S. 409 beginnend Zeile 5 v. u. bis S. 412, untermischt
mit späteren Aufzeichnungen aus anderen Handbüchern (17, Bc und 19, Be)
abgedruckt sind. Man wird wohl die Entstehungszeit dieser Aufzeichnungen mit
dem Beginn von \textsc{Gauss’} Untersuchungen zur Theorie der höheren Potenzreste,
also\footnote{Siehe die Tagebuchnotizen 130.—133. nebst den zugehörigen Bemerkungen, Werke X, 1, S. 565, 566 und \textsc{Bachmann}, Abh. 1 dieses Bandes, Art. 21, S. 52.} etwa auf 1805—1807, anzusetzen haben.\par
\HistoricalFoot{}{16\textsuperscript{*}}
\clearpage\SourceLeaf{209}\ArithmeticHead{124}{L}{L. SCHLESINGER, ÜBER GAUSS’ FUNKTIONENTHEORETISCHE ARBEITEN.}
In der Tat enthalten diese Notizen nicht nur eine Rekapitulation der
älteren (in \textit{Leiste} und Scheda Aa befindlichen) Potenzreihenentwicklungen für
den \(\sin\operatorname{lemn}\) und dessen Zähler und Nenner, sowie der Formeln für die ganzzahlige
und die komplexe Multiplikation der Zähler und Nenner von \(\sin\operatorname{lemn}\)
und \(\cos\operatorname{lemn}\), sondern am Schluss noch eine bemerkenswerte zahlentheoretische
Anwendung (diese stammt von S. 63 und 72 des Handbuchs 16), die als
Vorarbeit anzusehen sein dürfte für die in der Tagebuchnotiz 146. vom 9. Juli
1814 enthaltene »durch Induktion gemachte Beobachtung, die einen eleganten
Zusammenhang zwischen der Theorie der biquadratischen Reste und den
lemniskatischen Funktionen herstellt«. Diese »Beobachtung« besagt, dass, wenn
\(a+bi\) eine komplexe Primzahl bedeutet, für die \(a-1+bi\) durch \(2+2i\) teilbar
ist, die Anzahl der Lösungen der Kongruenz
\[
x^2+y^2+x^2y^2\equiv1\ (\bmod\ a+bi),
\]
\(x=\infty,\ y=\pm i,\ x=\pm i,\ y=\infty\) eingeschlossen, gleich der Norm \((a-1)^2+b^2\)
von \(a-1+bi\) sei. Nachdem \textsc{Dedekind} die Richtigkeit für alle \(a+bi\), deren
Norm unterhalb 100 liegt, verifiziert, und \textsc{Fricke} darauf hingewiesen hatte,
dass \(x^2+y^2+x^2y^2=0\) die Relation sei, die zwischen \(\operatorname{sl}\varphi\) und \(\operatorname{cl}\varphi=\operatorname{sl}\bigl(\frac{\varpi}{2}-\varphi\bigr)\)
besteht, hat G. \textsc{Herglotz} im Jahre 1921\footnote{Berichte der Sächsischen Akademie, math.-phys. Klasse 73, 1921, S. 271.} einen Beweis der \textsc{Gaussschen} Aussage
mit Hilfe von lemniskatischen Funktionen gegeben.

Dass auch für die oben (S. 102) wiedergegebene Aufzeichnung zur Theorie
der Modulfunktion aus der Scheda An, die ihrer Entstehungszeit nach hierher
gehört, ein Zusammenhang mit zahlentheoretischen Überlegungen vorzuliegen
scheint, haben wir bereits hervorgehoben.

Damit würden wohl die auf die Jahre 1801—1807 zu datierenden Untersuchungen
über Gegenstände der Analysis, soweit sich im Nachlass Spuren
von ihnen finden, erschöpft sein\footnote{Auf die aus dem Anfang des Jahres 1806 stammende erste Aufzeichnung über die Reihe \(F(\alpha,\beta,\gamma,x)\) kommen wir im Abschnitt VII. zurück.}; \textsc{Gauss} hat eben in diesen Jahren den
grössten Teil seiner Zeit astronomischen Arbeiten gewidmet.

Mit diesen astronomischen Arbeiten im Zusammenhang steht auch die
\textit{Theoria interpolationis methodo nova tractata}, die aus dem Nachlass Werke III,
S. 265—327 abgedruckt ist, und die wir hier darum erwähnen, weil eine erste
Redaktion dieser Abhandlung, die das Handbuch 18, Bd eröffnet, im Oktober\par
\clearpage\SourceLeaf{210}\ArithmeticHead{125}{R}{VIb. KORRESPONDENZ MIT SCHUMACHER 1808.}
\noindent 1805 verfasst ist\footnote{Vergl. auch die Angabe von \textsc{Schering}, Werke III, S. 328. Von der Handschrift ist in zwei Briefen \textsc{Schumachers} an \textsc{Gauss} die Rede; am 8. November 1808 schreibt er (\textit{Briefwechsel} I, S. 9): »Sie haben hier [d. h. in Altona] einen Kragen vergessen, auch hat meine Mutter noch Ihr Manuskript über Interpolation« und am 8. Juni 1816 (ebenda, S. 128): »Ich habe in der Hoffnung, Sie würden dadurch Veranlassung finden, Ihre Theorie der Interpolation, die ich handschriftlich habe, bekannt zu machen, die Preisfrage unserer Gesellschaft für 1817 so abfassen lassen .\ .\ .«.}. Dies wird auch durch die Tagebuchaufzeichnung 124.
vom November 1805, »Die Theorie der Interpolation haben wir weiter ausgebildet«
bestätigt, ferner aber dadurch, dass im Art. 41 (Werke III, S. 325), wo
als Beispiel die Mittelpunktsgleichung der Juno behandelt wird, derselbe Wert
der Exzentrizität \(0{,}254236\) auftritt, wie im Art. 1 des auf die Konvergenz
der Entwicklungskoeffizienten der Mittelpunktsgleichung bezüglichen Stücks
(Werke X, 1, S. 421). Dieser Wert gehört zu den V. Elementen dieses kleinen
Planeten (siehe die Bemerkungen ebenda, S. 443), die im Mai 1805 veröffentlicht
wurden (Werke VI, S. 462). Vielleicht besteht sogar ein tieferer Zusammenhang
zwischen beiden Untersuchungen; \textsc{Gauss} stellt nämlich in der
\textit{Theoria Interpolationis} Untersuchungen über den »Rest« der Reihenentwicklungen
an (siehe Art. 18 ff.), und man könnte dies als einen Versuch ansehen,
um Grenzen für den Fehler festzustellen, der bei der interpolatorischen Entwicklung
einer periodischen Funktion entsteht.
